% Zdroj: PDF priklady-jaro-2023.pdf, fyzická str. 5. Sada starší (neznačený formát 2023), otázka studijního zaměření. Zařazeno: S3.
\section{Naivní Bayesovský klasifikátor}
\szzotazka{jaro 2023}{16}{Naivní Bayesovský klasifikátor}

\noindent\textbf{Zadání.}\par
\begin{enumerate}
  \item Popište princip naivního Bayesovského klasifikátoru. Na základě čeho (vzorec) přiřazuje
    objekty do tříd?
  \item Popište význam jednotlivých pravděpodobností v Bayesově vzorci. V čem spočívá naivnost
    klasifikátoru?
  \item Použitím Bayesovského klasifikátoru určete, zda student půjde do školy, jestliže vstane
    pozdě, nemá oblíbenou přednášku a prší. Pravděpodobnost určete na základě následujících pozorování.
\end{enumerate}

\begin{center}
\begin{tabular}{ccccc}
\toprule
Den & Vstanu & Je oblíbená přednáška? & Počasí & Jdu do školy \\
\midrule
1  & Brzo     & Ne  & Slunce & Ne  \\
2  & Normálně & Ne  & Désť   & Ne  \\
3  & Brzo     & Ne  & Slunce & Ano \\
4  & Pozdě    & Ne  & Slunce & Ano \\
5  & Pozdě    & Ano & Slunce & Ano \\
6  & Pozdě    & Ano & Désť   & Ne  \\
7  & Brzo     & Ano & Désť   & Ano \\
8  & Normálně & Ne  & Slunce & Ne  \\
9  & Normálně & Ano & Slunce & Ano \\
10 & Pozdě    & Ano & Slunce & Ano \\
11 & Normálně & Ano & Désť   & Ano \\
12 & Brzo     & Ne  & Désť   & Ano \\
13 & Brzo     & Ano & Slunce & Ano \\
\bottomrule
\end{tabular}
\end{center}
{\footnotesize\noindent V~originálním PDF je hodnota psána „Désť``, správně „Déšť``; tabulka je přepsána věrně.\par}

\medskip
\noindent\textbf{Řešení.} \textit{(V~PDF bez náčrtu řešení; vlastní vzorové řešení, četnosti a skóre ověřeny numericky.)}\par
\begin{enumerate}[nosep]
  \item Klasifikátor přiřadí třídu maximalizující $P(c)\prod_j P(x_j\mid c)$ -- Bayesovo pravidlo s~\emph{naivním} předpokladem podmíněné nezávislosti příznaků. Pravděpodobnosti se odhadnou relativními četnostmi z~tabulky.
  \item $P(c)$ jsou apriorní pravděpodobnosti tříd, $P(x_j\mid c)$ věrohodnosti jednotlivých příznaků, jmenovatel $P(x_1,\dots,x_D)$ je normalizace (stejná pro všechny třídy, pro výběr maxima ji lze vynechat). \uv{Naivnost} je v~předpokladu, že příznaky jsou při dané třídě navzájem nezávislé.
  \item Z~13 dnů jde do školy $9\times$, nejde $4\times$: $P(\text{Ano})=\tfrac{9}{13}$, $P(\text{Ne})=\tfrac{4}{13}$. Pro dotaz (Vstanu${=}$Pozdě, oblíbená${=}$Ne, Počasí${=}$Déšť):
  \[ \text{Ano}:\ \tfrac{9}{13}\cdot\tfrac39\cdot\tfrac39\cdot\tfrac39=\tfrac{1}{39}\approx 0{,}026,\qquad
     \text{Ne}:\ \tfrac{4}{13}\cdot\tfrac14\cdot\tfrac34\cdot\tfrac24=\tfrac{3}{104}\approx 0{,}029. \]
  Větší je skóre třídy \textbf{Ne}, takže klasifikátor predikuje, že student \emph{do školy nepůjde} (po normalizaci $P(\text{Ne}\mid\text{data})\approx 53\,\%$).
\end{enumerate}
